\documentclass[11pt]{article}
\usepackage[margin=1in]{geometry}
\usepackage{booktabs}
\usepackage{amsmath,amssymb}
\usepackage{array}
\usepackage[T1]{fontenc}
\usepackage{lmodern}
\usepackage{microtype}
\usepackage{enumitem}
\usepackage{hyperref}
\hypersetup{colorlinks=true,linkcolor=black,urlcolor=blue}
\newcommand{\gate}{\texttt{g57}}
\title{Preimage cryptanalysis of the baseline mixed output:\\ efforts and conclusions}
\date{2026-10-26}
\author{}
\begin{document}
\maketitle

\section{Objective and setup}

The target is the declared mixed output of the baseline / as-is run of the TDP
pipeline: a reversible circuit \textbf{F} on 256 wires with \textbf{666{,}745 gates}
that computes a 64-bit permutation \textbf{C} on the honest input slice. The honest
input places $x$ on wires 0--63 and zeros elsewhere; the payload $C(x)$ appears on
wires 64--127 at the output.

\paragraph{Goal.} Find an input $x$ with $C(x)=0^{64}$ (a preimage of zero), using only
the declared output \textbf{F} and the attacker's own rewrites of it. The identity of
the internal $C$ values exposed by the earlier census is \emph{not} used --- it was
derived with knowledge of $C$ and is therefore off-limits. $C$ itself appears only in
clearly labelled diagnostics that measure what an attack could have seen; it never
drives an attack. Every candidate $x$ is verified against the real circuit.

Three equivalent forms of the same $C$ were available and used as targets:

\begin{center}
\begin{tabular}{lrl}
\toprule
circuit & gates & origin \\
\midrule
F & 666{,}745 & declared baseline/as-is mixed output \\
plateau & 355{,}990 & C-blind ANF resynthesis of F (30 rounds) \\
short-279k & 279{,}263 & fcompress shortening sweep (\texttt{fc4\_c5\_n64}), verified equivalent \\
\bottomrule
\end{tabular}
\end{center}

A preimage found on any of them is the same $x$, verifiable on all.

\section{Attacks attempted}

\subsection{Direct SAT inversion}
Each circuit was encoded into CNF (one variable per initial wire and per gate birth;
Toffoli gates as an AND auxiliary plus a native XOR clause for CryptoMiniSat), with the
non-input wires fixed to $0$ at entry and the payload fixed to $0^{64}$ at exit.
Instances on F and on the plateau ran with 8 threads each.
\textbf{Outcome:} no solution; both baselines ran more than 30 hours without returning.

\subsection{Structural reconstruction (vocabulary walks)}
Build the vocabulary of every function (values on 8{,}192 sampled inputs, up to
complement) born on any wire of the derived circuits, then walk \gate{} gates forward
from the inputs and backward from the payload, accepting a gate whenever the implied
value is present in the vocabulary. A complete chain would reconstruct $C$, hence $x$.
\textbf{Outcome:} no chain. The vocabulary covers 684 of $C$'s 1{,}000 gate births with
seven circuits and 735 with all 65 fair-game circuits, but mere membership does not
identify a genuine $C$ gate: the C-blind resyntheses and stage-3 remixes contain many
non-$C$ values of exactly the tested form. Forward walks were spurious from the first
gate; backward walks were genuine only with the smallest vocabulary and stalled after
20--40 gates. A lineage rule and a degree-2 span filter reduced false matches but did
not make the walk reliable.

\subsection{The replication landmark}
The raw sandwich computes $C(x)$ in the middle and copies it, so an unmixed circuit has
an interior cut where two disjoint 64-wire bundles carry the same functions; constraining
that bundle to zero would isolate $C$ in the first half. We tested, C-blind, whether this
survives, at three levels of generality.
\begin{itemize}[leftmargin=1.4em]
\item \textbf{Verbatim / up to NOT --- absent.} Over every cut of F and its three
  rewrites, the number of functions carried by two or more wires peaks at only 15--21 and
  only at the very start (depth $0.005$--$0.009$, the input-copy step); it is zero through
  the whole interior. F's 64 output functions are never co-present at an interior cut
  (all 64 only at the final gate). Stage-6 packing and ANF resynthesis rebuild the circuit
  toward the output, destroying the clean mid-circuit copy.
\item \textbf{Affine (degree-1) --- weak.} F peaks at only 10 of 64 output bits in the
  affine span of the live wires at depth $0.485$; the two remixes at 5. Insufficient to
  invert.
\item \textbf{Degree-2 --- the finding (Section~\ref{sec:finding}).}
\end{itemize}

\subsection{Exploiting the degree-2 structure}
Three sound ways to turn the degree-2 midpoint relations into an inversion were built
and run; all add constraints logically implied by $\text{payload}=0$ and all are
self-verifying.
\begin{itemize}[leftmargin=1.4em]
\item \textbf{Two-circuit linked SAT.} Encode the plateau and F over a shared
  64-variable input, both payloads zero, plus clauses equating internal variables whose
  functions coincide on samples. The two equivalent circuits share \textbf{42{,}771}
  identical internal functions, so a value forced cheaply in the small circuit
  propagates into the large one. Variants with 6{,}000 and 30{,}000 links ($\sim$1.7M
  variables).
\item \textbf{Hint-augmented SAT.} The full plateau with $\text{payload}=0$, plus the
  sparse degree-2 relations imposed at one cut (50 hints) and at two cuts (94 hints,
  covering 59 of 64 output bits) as redundant implied constraints.
\item \textbf{Prefix-only SAT.} Encode only the prefix up to the degree-2 cut, impose
  the relations there, and enumerate the prefix's satisfying inputs, checking each
  against the full circuit. On the plateau, 59 constraints across two cuts pin the input
  to about $2^{5}$ candidates over 178k of the 356k gates. On short-279k the earlier
  midpoint gives a shorter prefix (112k gates) but only 40 constraints.
\end{itemize}
\textbf{Outcome:} no candidate or preimage from any variant. The prefix solve --- the
best-shaped, smallest instance --- produced zero satisfying inputs after 16+ hours on
the 279k circuit and 24+ hours on the plateau; the linked and hint variants likewise
returned nothing.

\section{Key finding: the degree-2 midpoint replica}
\label{sec:finding}

The replication landmark survives the pipeline \textbf{not verbatim and not affinely,
but up to degree 2}. At each cut we measured how many of the 64 output functions are
GF(2) polynomials of degree at most 2 in the live wires, over enough samples (36{,}864)
that the ${\sim}33{,}000$ degree-2 monomials stay below the sample count (membership
false-positive $\approx 2^{-4000}$, so the counts are not coincidental). The count rises
sharply toward the middle and peaks per circuit:

\begin{center}
\begin{tabular}{lrcc}
\toprule
circuit & peak deg-2 bits & depth & profile at 0.40/0.45/0.48/0.50 \\
\midrule
plateau (356k) & 50 of 64 & 0.45 & 13 / 50 / 30 / 44 \\
F (666k) & 36 of 64 & 0.48 & 3 / 20 / 36 / 19 \\
short-279k & 40 of 64 & 0.40 & 40 / 37 / 26 / 29 \\
remix (1.38M) & 20 of 64 & 0.48 & 2 / 12 / 20 / 12 \\
remix2 (1.11M) & 16 of 64 & 0.50 & 3 / 8 / 8 / 16 \\
\bottomrule
\end{tabular}
\end{center}

Three properties make this a genuine, usable weakness rather than an artifact:
\begin{itemize}[leftmargin=1.4em]
\item \textbf{Sparse.} Despite the near-full-rank degree-2 span, the representations are
  compact: on the plateau at depth $0.45$ the 50 relations use a median of 10 monomials
  (max 50), and four output bits (17, 31, 41, 57) equal a single interior wire verbatim.
  On the 279k circuit the median is 6.
\item \textbf{Nested and growing.} On the plateau at depth $0.45$, by minimal replica: 4
  bits are a bare wire, 5 are affine, 50 are degree-2, 14 have none. The classes nest, so
  higher degree exposes genuinely new bits --- the affine-to-quadratic jump is almost
  entirely new (45 of the 50 need a product).
\item \textbf{Complementary across circuits.} The plateau and F expose largely different
  bits (overlap 29). Two plateau cuts cover 59 of 64; adding F's cut reaches
  \textbf{62 of 64}. Almost every output bit is a compact degree-2 function of interior
  wires at some midpoint cut.
\end{itemize}

An incidental observation: each step of ANF resynthesis \emph{increased} the degree-2
exposure ($666\text{k}\to 36$, $356\text{k}\to 50$), while the fcompress-shortened 279k
circuit peaks lower (40) and earlier (depth $0.40$) --- so the low-degree leak comes
specifically from ANF resynthesis, not from compression in general.

\section{Conclusions}

\begin{enumerate}[leftmargin=1.6em]
\item \textbf{No preimage was found.} Across three equivalent sizes (666k, 356k, 279k)
  and five sound SAT formulations (full, two-circuit linked, degree-2 hinted at one and
  two cuts, half-circuit with the bundle constrained), no solver produced $x$ with
  $C(x)=0^{64}$. The best-shaped instance returned no satisfying assignment at all.
\item \textbf{The structural leak is real and nearly complete.} The sandwich's
  mid-circuit replication of $C$ survives at degree 2: around the midpoint, almost every
  output bit (62 of 64 across two circuits) is a sparse degree-2 function of the interior
  wires --- a concrete, quantified weakness the verbatim and affine scans missed.
\item \textbf{The leak does not translate into a feasible inversion by these means.} The
  obstacle is the solver, not the structure: inverting a 100k--180k-gate reversible
  circuit is intractable for CryptoMiniSat here even with the degree-2 relations pinning
  the input and the uncompute half removed. Shrinking from 666k to 279k did not cross the
  threshold.
\item \textbf{Interpretation.} The degree-2 replica is a property of the
  \emph{representation} --- resynthesis expresses $C$'s output as low-degree functions of
  its own intermediates --- not a shortcut to $x$. Exploiting it still requires solving
  for an input consistent with the prefix, the hard part. A different engine (algebraic /
  Gr\"obner over the sparse quadratic system, or a dedicated reversible-circuit inverter)
  is the next thing to try; standard SAT is not enough.
\end{enumerate}

\section{Tools and artifacts}
All tools are on the server under \texttt{\textasciitilde/threemask\_cmp/py6/}; results
under \texttt{\textasciitilde/threemask\_cmp/run/preimage/}. Direct SAT:
\texttt{sat\_invert.py}. Vocabulary walks: \texttt{vocab.py}, \texttt{blind\_walk3.py},
\texttt{vocab\_diag*.py}. Replication scans: \texttt{payload\_cut.py},
\texttt{replication\_scan.py}, \texttt{affine\_replica\_scan.py},
\texttt{degree2\_replica\_scan.py}. Combinations: \texttt{combo\_extract.py},
\texttt{classify\_replicas.py}. Exploits: \texttt{combined\_link\_sat.py},
\texttt{plateau\_hint\_sat.py}, \texttt{plateau\_hint\_sat2.py},
\texttt{prefix\_quad\_sat.py}. The full chronological log is in
\texttt{docs/notes/stage2\_stage6\_campaign\_20261021.md}.

\end{document}
